\chapter{Cauchy's Theorem}
\label{ch:iv05}

Cauchy's theorem is the topological heart of elementary complex analysis: holomorphic functions have zero integral around contractible closed contours. This creates path independence, primitives, and deformation invariance.

\section*{Learning goals}
After this chapter, the reader should be able to:
\begin{itemize}
\item verify the hypotheses of Cauchy's theorem on simply connected or triangulated domains;
\item deform contours through holomorphic regions without crossing singularities;
\item prove vanishing of closed contour integrals from holomorphicity and topology;
\item identify when multiply connected geometry prevents a direct use of Cauchy's theorem;
\item construct primitives from path integrals when closed integrals vanish;
\item analyze counterexamples showing why a puncture or singularity changes the conclusion;
\end{itemize}
\section*{Conceptual roadmap}
\[
\boxed{\text{local holomorphic structure}\;\longrightarrow\;\text{integral or series representation}\;\longrightarrow\;\text{global consequence}.}
\]

\section{Triangle theorem}
The theorem can first be proved on triangles, where subdivision and local linearization isolate the holomorphic error.

\section{Polygonal domains}
Triangulation extends the result to polygonal contours.

% BEGIN VOL04-EXPANSION IV05-example-01
\begin{example}[Triangle cancellation]\label{ex:iv05-expand-01}
Suppose \(f\) is holomorphic on a neighborhood of a triangle and its interior. Cauchy's theorem gives
\[
\int_{\partial T}f(z)\,dz=0.
\]
For a polynomial such as \(f(z)=z^3+2z\), this can also be checked from the primitive \(z^4/4+z^2\). The theorem is stronger because it applies even when a primitive has not yet been constructed explicitly.
\end{example}
% END VOL04-EXPANSION IV05-example-01
\section{Simply connected domains}
On simply connected regions, every closed contour is deformable to a point and holomorphic integrals vanish.

\section{Homotopy invariance}
Contour integrals of holomorphic functions are invariant under suitable fixed-endpoint homotopies.

\section{Existence of primitives}
Zero integrals on closed contours imply primitives.

% BEGIN VOL04-EXPANSION IV05-example-02
\begin{example}[Deforming a contour without crossing singularities]\label{ex:iv05-expand-02}
Let \(f(z)=1/(z-3)\), and let \(\gamma_0\) and \(\gamma_1\) be two closed curves lying in the disk \(|z|<2\). Since the pole at \(3\) lies outside the disk and the curves can be deformed into one another there, Cauchy's theorem gives equal integrals. In fact each integral is zero because \(f\) has a primitive on the simply connected disk.
\end{example}
% END VOL04-EXPANSION IV05-example-02
\section{Goursat refinement}
Continuity of the derivative is not required; complex differentiability alone suffices.

\section{Multiply connected warning}
Loops around holes can carry nonzero integrals even when the function is holomorphic on the punctured domain.

\section{Topological content}
Cauchy's theorem converts local holomorphicity plus global domain topology into global integral identities.

% BEGIN VOL04-EXPANSION IV05-example-03
\begin{example}[The puncture blocks Cauchy's theorem]\label{ex:iv05-expand-03}
On \(\mathbb C\setminus\{0\}\), the function \(1/z\) is holomorphic, yet
\[
\int_{|z|=1}\frac{dz}{z}=2\pi i.
\]
There is no contradiction: the unit circle does not bound a region contained in the domain. The missing point is a topological obstruction, and the example explains why domain hypotheses matter as much as pointwise holomorphicity.
\end{example}
% END VOL04-EXPANSION IV05-example-03
\section{Core structural results}

\begin{theorem}[Cauchy--Goursat theorem on triangles]\label{thm:iv05-01}
If $f$ is holomorphic on a neighborhood of a triangle and its interior, then the integral around the triangle is zero.
\begin{proof}
Repeatedly subdivide into four similar triangles and choose one carrying at least a quarter of the integral magnitude. Diameters shrink to a point where differentiability gives a linear approximation whose constant and linear parts integrate to zero; the remainder estimate forces the original integral to vanish.
\end{proof}
\end{theorem}

\begin{theorem}[Primitive on simply connected domains]\label{thm:iv05-02}
Every holomorphic function on a simply connected domain has a primitive.
\begin{proof}
Cauchy's theorem makes integrals around closed contours vanish, so path integration from a base point is well-defined and differentiates back to the integrand.
\end{proof}
\end{theorem}

\begin{theorem}[Homotopy invariance]\label{thm:iv05-03}
If two closed contours are homotopic through contours inside a domain where $f$ is holomorphic, their integrals agree.
\begin{proof}
Partition the homotopy strip into small cells whose boundary integrals vanish by local Cauchy theory; interior edges cancel, leaving the difference of the two boundary contour integrals.
\end{proof}
\end{theorem}

\section{Worked examples}

\begin{example}[Triangle theorem diagnostic]\label{ex:iv05-worked-01}
Give a rigorous worked analysis of Triangle theorem. State the relevant holomorphic, contour, series, or singularity hypothesis and derive the central conclusion. The theorem can first be proved on triangles, where subdivision and local linearization isolate the holomorphic error. The decisive step is to use the complex-analytic structure globally rather than reason from real-variable intuition alone.
\end{example}

\begin{example}[Polygonal domains diagnostic]\label{ex:iv05-worked-02}
Give a rigorous worked analysis of Polygonal domains. State the relevant holomorphic, contour, series, or singularity hypothesis and derive the central conclusion. Triangulation extends the result to polygonal contours. The decisive step is to use the complex-analytic structure globally rather than reason from real-variable intuition alone.
\end{example}

\begin{example}[Simply connected domains diagnostic]\label{ex:iv05-worked-03}
Give a rigorous worked analysis of Simply connected domains. State the relevant holomorphic, contour, series, or singularity hypothesis and derive the central conclusion. On simply connected regions, every closed contour is deformable to a point and holomorphic integrals vanish. The decisive step is to use the complex-analytic structure globally rather than reason from real-variable intuition alone.
\end{example}

\begin{example}[Homotopy invariance diagnostic]\label{ex:iv05-worked-04}
Give a rigorous worked analysis of Homotopy invariance. State the relevant holomorphic, contour, series, or singularity hypothesis and derive the central conclusion. Contour integrals of holomorphic functions are invariant under suitable fixed-endpoint homotopies. The decisive step is to use the complex-analytic structure globally rather than reason from real-variable intuition alone.
\end{example}

\section{Solved dossiers}
The dossiers are canonical solved problems. The provenance ledger separately records which retained corpus rules guided each topic and which dossiers were newly designed.

\begin{problem}[Triangle theorem diagnostic]\label{prob:iv05-dossier-01}
Give a rigorous worked analysis of Triangle theorem. State the relevant holomorphic, contour, series, or singularity hypothesis and derive the central conclusion.
\end{problem}
\begin{solution}
The theorem can first be proved on triangles, where subdivision and local linearization isolate the holomorphic error. The decisive step is to use the complex-analytic structure globally rather than reason from real-variable intuition alone.
\end{solution}

\begin{problem}[Polygonal domains diagnostic]\label{prob:iv05-dossier-02}
Give a rigorous worked analysis of Polygonal domains. State the relevant holomorphic, contour, series, or singularity hypothesis and derive the central conclusion.
\end{problem}
\begin{solution}
Triangulation extends the result to polygonal contours. The decisive step is to use the complex-analytic structure globally rather than reason from real-variable intuition alone.
\end{solution}

\begin{problem}[Simply connected domains diagnostic]\label{prob:iv05-dossier-03}
Give a rigorous worked analysis of Simply connected domains. State the relevant holomorphic, contour, series, or singularity hypothesis and derive the central conclusion.
\end{problem}
\begin{solution}
On simply connected regions, every closed contour is deformable to a point and holomorphic integrals vanish. The decisive step is to use the complex-analytic structure globally rather than reason from real-variable intuition alone.
\end{solution}

\begin{problem}[Homotopy invariance diagnostic]\label{prob:iv05-dossier-04}
Give a rigorous worked analysis of Homotopy invariance. State the relevant holomorphic, contour, series, or singularity hypothesis and derive the central conclusion.
\end{problem}
\begin{solution}
Contour integrals of holomorphic functions are invariant under suitable fixed-endpoint homotopies. The decisive step is to use the complex-analytic structure globally rather than reason from real-variable intuition alone.
\end{solution}

\begin{problem}[Existence of primitives diagnostic]\label{prob:iv05-dossier-05}
Give a rigorous worked analysis of Existence of primitives. State the relevant holomorphic, contour, series, or singularity hypothesis and derive the central conclusion.
\end{problem}
\begin{solution}
Zero integrals on closed contours imply primitives. The decisive step is to use the complex-analytic structure globally rather than reason from real-variable intuition alone.
\end{solution}

\begin{problem}[Goursat refinement diagnostic]\label{prob:iv05-dossier-06}
Give a rigorous worked analysis of Goursat refinement. State the relevant holomorphic, contour, series, or singularity hypothesis and derive the central conclusion.
\end{problem}
\begin{solution}
Continuity of the derivative is not required; complex differentiability alone suffices. The decisive step is to use the complex-analytic structure globally rather than reason from real-variable intuition alone.
\end{solution}

\begin{problem}[Multiply connected warning diagnostic]\label{prob:iv05-dossier-07}
Give a rigorous worked analysis of Multiply connected warning. State the relevant holomorphic, contour, series, or singularity hypothesis and derive the central conclusion.
\end{problem}
\begin{solution}
Loops around holes can carry nonzero integrals even when the function is holomorphic on the punctured domain. The decisive step is to use the complex-analytic structure globally rather than reason from real-variable intuition alone.
\end{solution}

\begin{problem}[Topological content diagnostic]\label{prob:iv05-dossier-08}
Give a rigorous worked analysis of Topological content. State the relevant holomorphic, contour, series, or singularity hypothesis and derive the central conclusion.
\end{problem}
\begin{solution}
Cauchy's theorem converts local holomorphicity plus global domain topology into global integral identities. The decisive step is to use the complex-analytic structure globally rather than reason from real-variable intuition alone.
\end{solution}

\begin{problem}[Cauchy--Goursat theorem on triangles proof dossier]\label{prob:iv05-dossier-09}
Prove the following theorem-level statement and identify the essential hypothesis: If $f$ is holomorphic on a neighborhood of a triangle and its interior, then the integral around the triangle is zero.
\end{problem}
\begin{solution}
Repeatedly subdivide into four similar triangles and choose one carrying at least a quarter of the integral magnitude. Diameters shrink to a point where differentiability gives a linear approximation whose constant and linear parts integrate to zero; the remainder estimate forces the original integral to vanish. This identifies the mechanism that makes the complex-analytic conclusion stronger than its real-variable analogue.
\end{solution}

\begin{problem}[Primitive on simply connected domains proof dossier]\label{prob:iv05-dossier-10}
Prove the following theorem-level statement and identify the essential hypothesis: Every holomorphic function on a simply connected domain has a primitive.
\end{problem}
\begin{solution}
Cauchy's theorem makes integrals around closed contours vanish, so path integration from a base point is well-defined and differentiates back to the integrand. This identifies the mechanism that makes the complex-analytic conclusion stronger than its real-variable analogue.
\end{solution}

\begin{problem}[Homotopy invariance proof dossier]\label{prob:iv05-dossier-11}
Prove the following theorem-level statement and identify the essential hypothesis: If two closed contours are homotopic through contours inside a domain where $f$ is holomorphic, their integrals agree.
\end{problem}
\begin{solution}
Partition the homotopy strip into small cells whose boundary integrals vanish by local Cauchy theory; interior edges cancel, leaving the difference of the two boundary contour integrals. This identifies the mechanism that makes the complex-analytic conclusion stronger than its real-variable analogue.
\end{solution}

\begin{problem}[Chapter synthesis]\label{prob:iv05-dossier-12}
Connect the definitions and principal theorems of this chapter into one reusable complex-analysis strategy.
\end{problem}
\begin{solution}
Cauchy's theorem is the topological heart of elementary complex analysis: holomorphic functions have zero integral around contractible closed contours. This creates path independence, primitives, and deformation invariance. The reusable strategy is to identify the analytic domain, choose the correct local expansion or contour representation, apply the structural theorem, and then audit singularities and topology.
\end{solution}

\section{Exercises with complete solutions}

\begin{exercise}\label{exr:iv05-01}
State and apply the central principle behind Triangle theorem. Give one concrete consequence.
\end{exercise}
\begin{hint}
Check holomorphicity on an open set containing the contour and its interior before invoking Cauchy's theorem.
\end{hint}
\begin{solution}
The theorem can first be proved on triangles, where subdivision and local linearization isolate the holomorphic error. The same principle can then be reused in later contour, residue, conformal, or Riemann-surface arguments.
\end{solution}

\begin{exercise}\label{exr:iv05-02}
State and apply the central principle behind Polygonal domains. Give one concrete consequence.
\end{exercise}
\begin{hint}
During contour deformation, verify the swept region is free of singularities; homotopy through a puncture is not allowed.
\end{hint}
\begin{solution}
Triangulation extends the result to polygonal contours. The same principle can then be reused in later contour, residue, conformal, or Riemann-surface arguments.
\end{solution}

\begin{exercise}\label{exr:iv05-03}
State and apply the central principle behind Simply connected domains. Give one concrete consequence.
\end{exercise}
\begin{hint}
Triangulate or decompose the region so interior edge integrals cancel pairwise.
\end{hint}
\begin{solution}
On simply connected regions, every closed contour is deformable to a point and holomorphic integrals vanish. The same principle can then be reused in later contour, residue, conformal, or Riemann-surface arguments.
\end{solution}

\begin{exercise}\label{exr:iv05-04}
State and apply the central principle behind Homotopy invariance. Give one concrete consequence.
\end{exercise}
\begin{hint}
In a multiply connected region, inspect each hole; closed contour integrals need not vanish around omitted points.
\end{hint}
\begin{solution}
Contour integrals of holomorphic functions are invariant under suitable fixed-endpoint homotopies. The same principle can then be reused in later contour, residue, conformal, or Riemann-surface arguments.
\end{solution}

\begin{exercise}\label{exr:iv05-05}
State and apply the central principle behind Existence of primitives. Give one concrete consequence.
\end{exercise}
\begin{hint}
Define a candidate primitive by a path integral from a base point and use vanishing of closed integrals to prove path independence.
\end{hint}
\begin{solution}
Zero integrals on closed contours imply primitives. The same principle can then be reused in later contour, residue, conformal, or Riemann-surface arguments.
\end{solution}

\begin{exercise}\label{exr:iv05-06}
State and apply the central principle behind Goursat refinement. Give one concrete consequence.
\end{exercise}
\begin{hint}
For a punctured-domain counterexample, integrate \(1/z\) around a loop with nonzero winding number.
\end{hint}
\begin{solution}
Continuity of the derivative is not required; complex differentiability alone suffices. The same principle can then be reused in later contour, residue, conformal, or Riemann-surface arguments.
\end{solution}

\begin{exercise}\label{exr:iv05-07}
State and apply the central principle behind Multiply connected warning. Give one concrete consequence.
\end{exercise}
\begin{hint}
Use local primitives on small disks and patch them only when the overlap integrals are consistent.
\end{hint}
\begin{solution}
Loops around holes can carry nonzero integrals even when the function is holomorphic on the punctured domain. The same principle can then be reused in later contour, residue, conformal, or Riemann-surface arguments.
\end{solution}

\begin{exercise}\label{exr:iv05-08}
State and apply the central principle behind Topological content. Give one concrete consequence.
\end{exercise}
\begin{hint}
Separate topology from regularity: holomorphicity alone does not guarantee every closed integral vanishes on a non-simply-connected domain.
\end{hint}
\begin{solution}
Cauchy's theorem converts local holomorphicity plus global domain topology into global integral identities. The same principle can then be reused in later contour, residue, conformal, or Riemann-surface arguments.
\end{solution}

\section*{Chapter summary}
The chapter now forms part of the canonical complex-analysis chain from holomorphicity through residues, global continuation, and Riemann surfaces.
% BEGIN VOL04-EXPANSION IV05-exercises-01
\section{Graded supplementary exercises}
These exercises extend the protected chapter with a balanced set of computations, proofs, hypothesis tests, applications, and challenges.

\subsection*{Standard computations}

\begin{exercise}[Polynomial loop]\label{exr:iv05-expand-01}
Evaluate \(\int_\gamma (z^4+z)\,dz\) over any closed contour.
\end{exercise}
\begin{hint}
The integrand has an entire primitive.
\end{hint}
\begin{solution}
The value is zero.
\end{solution}

\begin{exercise}[Disk holomorphic]\label{exr:iv05-expand-02}
Let \(f(z)=1/(z-5)\). Evaluate its integral around \(|z|=2\).
\end{exercise}
\begin{hint}
The function is holomorphic on the disk bounded by the contour.
\end{hint}
\begin{solution}
Cauchy's theorem gives zero.
\end{solution}

\begin{exercise}[Rectangle]\label{exr:iv05-expand-03}
Evaluate the integral of \(e^z\) around the boundary of a rectangle.
\end{exercise}
\begin{hint}
The exponential is entire.
\end{hint}
\begin{solution}
The integral is zero.
\end{solution}

\begin{exercise}[Two paths]\label{exr:iv05-expand-04}
If \(f\) is holomorphic on a simply connected domain, compare integrals along two paths with the same endpoints.
\end{exercise}
\begin{hint}
Join one path to the reverse of the other.
\end{hint}
\begin{solution}
The resulting closed contour has integral zero, so the two path integrals agree.
\end{solution}

\begin{exercise}[Primitive consequence]\label{exr:iv05-expand-05}
What is \(\int_\gamma \cos z\,dz\) from \(0\) to \(\pi\)?
\end{exercise}
\begin{hint}
Use the primitive \(\sin z\).
\end{hint}
\begin{solution}
The value is \(\sin\pi-\sin0=0\).
\end{solution}

\subsection*{Proofs}

\begin{exercise}[Path independence from closed loops]\label{exr:iv05-expand-06}
Prove that vanishing of every closed contour integral implies path independence.
\end{exercise}
\begin{hint}
Concatenate one path with the reverse of the other.
\end{hint}
\begin{solution}
The closed integral is the first path integral minus the second. If it vanishes, the two are equal.
\end{solution}

\begin{exercise}[Primitive from path independence]\label{exr:iv05-expand-07}
On a connected domain with path independent integrals, construct a primitive of \(f\).
\end{exercise}
\begin{hint}
Fix a base point and integrate to \(z\).
\end{hint}
\begin{solution}
Define \(F(z)=\int_{z_*}^{z}f(w)dw\). Path independence makes this well defined. A short final segment and continuity of \(f\) show \(F'(z)=f(z)\).
\end{solution}

\begin{exercise}[Simply connected implication]\label{exr:iv05-expand-08}
Explain why Cauchy's theorem yields primitives on a simply connected domain.
\end{exercise}
\begin{hint}
Closed curves can be contracted without leaving the domain.
\end{hint}
\begin{solution}
Cauchy's theorem makes every closed contour integral vanish; path independence follows, and the base point integral construction produces a primitive.
\end{solution}

\begin{exercise}[Contour deformation]\label{exr:iv05-expand-09}
Prove that two homologous boundary curves have equal integrals when \(f\) is holomorphic in the region between them.
\end{exercise}
\begin{hint}
Orient the boundary of the region between the curves.
\end{hint}
\begin{solution}
The boundary consists of one curve and the reverse of the other. Cauchy's theorem makes the total integral zero, hence the two original integrals are equal.
\end{solution}

\subsection*{Counterexamples and hypothesis tests}

\begin{exercise}[Punctured disk]\label{exr:iv05-expand-10}
Why can Cauchy's theorem not be used to conclude \(\int_{|z|=1}dz/z=0\)?
\end{exercise}
\begin{hint}
Inspect the interior of the contour.
\end{hint}
\begin{solution}
The interior contains the missing point \(0\), where the integrand is not holomorphic. The required filled region is not contained in the domain.
\end{solution}

\begin{exercise}[Pole on the boundary]\label{exr:iv05-expand-11}
Can the standard Cauchy theorem be applied when the integrand has a pole on the contour?
\end{exercise}
\begin{hint}
Holomorphicity is required on a neighborhood of the contour and its interior.
\end{hint}
\begin{solution}
No. The contour integral may even be undefined in the ordinary sense, so the hypothesis fails before the theorem is invoked.
\end{solution}

\begin{exercise}[Disconnected domain]\label{exr:iv05-expand-12}
Does simple connectedness make sense without connectedness?
\end{exercise}
\begin{hint}
Check the definition of a simply connected domain.
\end{hint}
\begin{solution}
A domain is connected and open by convention. Statements about simple connectedness presuppose that connected setting.
\end{solution}

\subsection*{Applications and investigations}

\begin{exercise}[Contour simplification]\label{exr:iv05-expand-13}
Explain how Cauchy's theorem can simplify a difficult path to an easier one.
\end{exercise}
\begin{hint}
Deform the path through a region free of singularities.
\end{hint}
\begin{solution}
The integral is unchanged under such a deformation. One can therefore replace a complicated contour by a geometrically convenient representative of the same deformation class.
\end{solution}

\begin{exercise}[Detecting holes]\label{exr:iv05-expand-14}
How can a nonzero integral of a holomorphic function around a closed contour reveal domain topology?
\end{exercise}
\begin{hint}
Compare with the simply connected conclusion.
\end{hint}
\begin{solution}
A nonzero closed integral shows that no primitive exists on the whole domain and that the contour cannot be contracted through a region where the integrand remains holomorphic.
\end{solution}

\subsection*{Challenge problems}

\begin{exercise}[Morera converse]\label{exr:iv05-expand-15}
State the idea behind the converse principle that vanishing triangle integrals can imply holomorphicity.
\end{exercise}
\begin{hint}
Use triangle integrals to construct local primitives.
\end{hint}
\begin{solution}
If a continuous function has zero integral around every triangle, one constructs a path independent local integral and hence a local primitive. Differentiating that primitive recovers the function, proving holomorphicity.
\end{solution}

\begin{exercise}[Annulus comparison]\label{exr:iv05-expand-16}
Two circles \(|z|=1\) and \(|z|=2\) lie in an annulus where \(f\) is holomorphic. Show their positively oriented integrals are equal.
\end{exercise}
\begin{hint}
Apply Cauchy's theorem to the region between them with boundary orientation.
\end{hint}
\begin{solution}
The boundary integral is the outer circle integral minus the inner circle integral. It vanishes, so the two are equal.
\end{solution}

% END VOL04-EXPANSION IV05-exercises-01
